\section{The actual slow-band Gram asymptotic}
\label{sec:slow-band-gram}

The preceding counterpath does not address the repaired regime
$Y/X\to0$.  This section proves the actual global asymptotic
in that regime at polynomial weight zero.  The proof retains
the weighted Gram projection and estimates its trace correction.
It requires no generic choice of strip edge and no separation
of zeros from that edge.

\subsection{An abstract trace-class Hankel tail}

\paragraph{MA-SB-001 (uniform rectangles from the actual model-space lift).}
Let $B$ be an inner function on the right half-plane.
Use the Laplace isometry
$\mathcal L:L^2(0,\infty)\to H^2(\operatorname{Re}z>0)$
with boundary norm $(2\pi)^{-1}\int|F(it)|^2dt$.
On $\mathcal K_+=L^2(0,\infty)$ put
\[
 V=\mathcal L^{-1}M_B\mathcal L,\qquad
 P_K=1-VV^* .
 \tag{MA-SB.1}
\]
Thus $V$ is an isometry and $P_K$ is the actual orthogonal
model-space projection in time coordinates.
Let $k\in C_c^\infty((0,\infty))$ and let
$C_k^+$ be its causal convolution on $\mathcal K_+$.
Assume the \emph{actual ambient lift}
\[
                         C_k^+P_K\in\mathcal S_1(\mathcal K_+).
 \tag{MA-SB.2}
\]
This is a hypothesis about an actual operator, stronger
than merely assigning a summable scalar spectrum.
For the zeta inner function used below, MA-GP-009 proves it
without an edge-separation assumption.

Let $\mathcal K=L^2(\mathbb R)$, and let $P_+$ and $P_-$
be its positive- and negative-half-line projections.
Multiplication by the boundary function $B(it)$ defines
a full-line unitary convolution operator $\mathcal V$.
Define
\[
 \begin{aligned}
 H_B&=P_+\mathcal V|_{\mathcal K_-},&
 V_-&=P_-\mathcal V|_{\mathcal K_-},\\
 H_k&=P_+C_k|_{\mathcal K_-},&
 H_{Bk}&=P_+C_k\mathcal V|_{\mathcal K_-}.
 \end{aligned}
 \tag{MA-SB.3}
\]
The symbol of the last operator is
$B(it)\int_0^\infty k(r)e^{-itr}dr$.
Then
\[
 \begin{gathered}
 H_BH_B^*=P_K,\qquad H_B=P_KH_B,\\
 H_{Bk}=C_k^+P_KH_B+H_kV_-\in
                  \mathcal S_1(\mathcal K_-,\mathcal K_+).
 \end{gathered}
 \tag{MA-SB.4}
\]
For $R>L\ge0$ the actual tail rectangle satisfies
\[
 \left\|P_{[R,\infty)}VC_k^+P_{[0,L]}\right\|_1
 \le
 \left\|P_{[R-L,\infty)}H_{Bk}\right\|_1.
 \tag{MA-SB.5}
\]
The right side tends to zero as $R-L\to\infty$,
uniformly in $L$.  The same bound holds if the output
projection on the left is replaced by any multiplication
operator of norm at most one supported in $[R,\infty)$.

\begin{proof}
The Hardy invariance corresponding to causality gives
$P_-\mathcal VP_+=0$.  Relative to
$\mathcal K_+\oplus\mathcal K_-$, its block form is
\[
                      \mathcal V=
                       \begin{pmatrix}V&H_B\\0&V_-\end{pmatrix}.
 \tag{MA-SB.6}
\]
The first diagonal block of $\mathcal V\mathcal V^*=1$
gives $VV^*+H_BH_B^*=1$, proving the first identity
in (MA-SB.4).
Since $P_K$ is an orthogonal projection,
$(1-P_K)H_BH_B^*(1-P_K)=0$, and hence
$(1-P_K)H_B=0$.
All the displayed blocks have norm at most one.

After identifying a negative variable with its opposite,
the kernel of $H_k$ is $k(x+y)$ for $x,y>0$.
It is the restriction of a smooth kernel supported on
a fixed compact square: choose smooth factors equal to
one on $[0,b]$, where $\operatorname{supp}k\subset(0,b)$.
Multiplying $k(x+y)$ by those factors does not change it
on the positive quadrant.
A double Fourier expansion on a larger square is
absolutely summable after repeated integration by parts.
It gives a rank-one trace-class representation; restricting
to the positive quadrant preserves trace class.
Thus $H_k\in\mathcal S_1$.

Inserting $P_++P_-$ between $C_k$ and $\mathcal V$ gives
\[
                      H_{Bk}=C_k^+H_B+H_kV_- .
 \tag{MA-SB.7}
\]
The preceding identity $H_B=P_KH_B$ and hypothesis
(MA-SB.2) prove the trace-class assertion in (MA-SB.4).
No differentiability of $B(it)$ is used.
In particular zeros may approach the boundary arbitrarily
closely at large heights.

The full-line convolution operators $C_k$ and
$\mathcal V$ commute.  On positive-half-line inputs their
product agrees with $VC_k^+$.
Extend the operator in (MA-SB.5) by zero to $\mathcal K$
and translate both input and output coordinates by $L$.
Translation commutes with this full-line convolution.
The resulting operator is exactly
\[
        P_{[R-L,\infty)}\,C_k\mathcal V\,P_{[-L,0]}
       =P_{[R-L,\infty)}\,H_{Bk}\,P_{[-L,0]} .
 \tag{MA-SB.8}
\]
Translations are unitary and compression is contractive
for the trace norm, proving (MA-SB.5).
Finally $P_{[r,\infty)}\to0$ strongly on $\mathcal K_+$
as $r\to\infty$.
Left multiplication of a fixed trace-class operator
by these uniformly bounded projections converges in
trace norm: this holds first for rank-one operators
and then by finite-rank approximation.
The multiplication-profile version follows by factoring
its norm-one multiplier through the same tail projection.
\end{proof}

\subsection{The weighted Gram correction on the positive tail}

Fix $0<\alpha<1/2$, polynomial weight $s=0$, and
$g\in C_c^\infty((0,\infty))$.
Put
\[
 \begin{gathered}
 H=L^2(\mathbb R,e^{2\alpha|x|}\,dx),\qquad
 k(t)=e^{\alpha t}g(t),\\
 \sigma_+=\tfrac12+\alpha,\qquad
 \tau_+(g)=\operatorname{Tr}_{\mathcal K_+}(C_k^+P_K),
 \end{gathered}
 \tag{MA-SB.9}
\]
where $B=B_\alpha$, $V$ and $P_K$ are the actual
half-plane objects of MA-GP-003/009.
The same fixed $B$ and $k$ are used for every band $Y$.
MA-GP-009, including its endpoint term at zero, gives
(MA-SB.2) and identifies this trace with the actual
finite-level quotient trace.

For $Y\ge1$ write $L=\log Y$ and use $r=x+L\ge0$ on
the actual support half-line $x\ge-\log Y$.
Its metric is
\[
 R_L(r)=
 \begin{cases}
 e^{4\alpha(L-r)},&0\le r\le L,\\
 1,&r\ge L.
 \end{cases}
 \tag{MA-SB.10}
\]
The map $f(x)\mapsto e^{\alpha x}f(x)$, followed by this
translation of coordinates, is unitary from that actual
weighted half-line into $L^2(R_L(r)\,dr)$.
Let
\[
 \mathfrak G_L=V^*M_{R_L}V,\qquad
 D_L=\mathfrak G_L^{-1}-1,\qquad
 Q_L^w=V\mathfrak G_L^{-1}V^*M_{R_L}.
 \tag{MA-SB.11}
\]
The last operator is the actual orthogonal projection
for the $R_L$ metric, as in MA-GP-004.
The identity map from $L^2(R_Ldr)$ to $\mathcal K_+$
is a bounded isomorphism for each fixed $L$.
It preserves traces under conjugation, but is not
being asserted uniformly isometric as $L$ grows.

\paragraph{MA-SB-002 (uniform actual slow-band trace at every strip edge).}
Define the two fixed trace-norm tails
\[
 \begin{aligned}
 \omega_k(r)&=\|P_{[r,\infty)}C_k^+P_K\|_1,\\
 \nu_k(r)&=\|P_{[r,\infty)}H_{Bk}\|_1 .
 \end{aligned}
 \tag{MA-SB.12}
\]
Both tend to zero as $r\to\infty$, by
MA-GP-009 and MA-SB-001.
Let $0\le\phi\le1$ be a compactly supported
multiplication profile on $[0,\infty)$ equal to one
on $[0,R]$, where $R>L$.
Then the actual weighted trace satisfies
\[
 \left|
 \operatorname{Tr}_{L^2(R_Ldr)}(M_\phi C_k^+Q_L^w)
                     +\tau_+(g)\right|
            \le\omega_k(R)+\nu_k(R-L).
 \tag{MA-SB.13}
\]
In particular let $\rho$ be any smooth source cutoff
of Definition~9.3.3.1, supported on $[1/X,X]$ and
equal to one on $[1/X+\delta,X-\delta]$.
If
\[
           Y\ge1,\qquad X\ge2Y,\qquad
           0<\delta\le1/(4X),
 \tag{MA-SB.14}
\]
then, with $Q_Y^{\rm glob}$ the actual global PW
projection,
\[
 \begin{aligned}
 &\left|\operatorname{Tr}_H(M_\rho C_gQ_Y^{\rm glob})
                    +\tau_+(g)\right|\\
 &\qquad\le
 \omega_k(\log Y+\log(X-\delta))
                 +\nu_k(\log(X-\delta))\\
 &\qquad\le
 \omega_k(\log(X/2))+\nu_k(\log(X/2)).
 \end{aligned}
 \tag{MA-SB.15}
\]
As before, $M_\rho$ denotes $\rho(e^x)$ in logarithmic
coordinates.
Thus along every path with $X,Y\to\infty$,
$Y/X\to0$ and the source's thin smoothing condition,
\[
           \operatorname{Tr}_H(M_\rho C_gQ_Y^{\rm glob})
                      \longrightarrow-\tau_+(g).
 \tag{MA-SB.16}
\]
This limit is uniform in $1\le Y\le X/2$ in the
explicit sense of (MA-SB.15).
No estimate of the rate at which the two fixed
tails themselves decay is assumed.
By MA-GP-010 the same conclusion holds for the
actual finite-$S$ invariant global PW projector
$\widehat Q_{S,Y}$, with $S'=S$.

\begin{proof}
\emph{The Gram correction has uniformly bounded support and norm.}
Put $P_L=P_{[0,L]}$ on $\mathcal K_+$.
Since $V$ is causal,
$P_LV=P_LVP_L$ and $V^*P_L=P_LV^*P_L$.
The difference $R_L-1$ is supported on $[0,L]$.
Consequently
\[
       \mathfrak G_L-1=P_L(\mathfrak G_L-1)P_L,\qquad
       D_L=P_LD_LP_L,\qquad -1\le D_L\le0 .
 \tag{MA-SB.17}
\]
Indeed $\mathfrak G_L\ge1$ and it is the identity
on $P_L^\perp$; its inverse has the same block
decomposition and lies between zero and one.
In particular $\|D_L\|\le1$, independently of $L$.
This is a support statement for the inverse Gram
correction, not a uniform bound for the whole
projection difference with its right metric factor.

\emph{All cyclic products are trace class before cancellation.}
In these half-line coordinates $M_\phi C_k^+$ is
trace class.  Its kernel is the compact output
cutoff times $k(r-s)$; causality and compact
support of $k$ also confine its input variable
to a fixed compact interval.
A smooth compact-square extension followed by
bounded coordinate multipliers proves trace
class, just as in MA-GP-008.
This proof applies also to arbitrary bounded compact profiles.
Its trace is $k(0)\int\phi=0$.
The diagonal assertion follows first for smooth
profiles by the same nuclear kernel expansion
and then for bounded compact profiles by bounded smooth
approximation, strong multiplier convergence and trace-norm convergence.
Similarly $P_LC_k^+$ is trace class.
Thus $D_LC_k^+$ is trace class by (MA-SB.17).

For this fixed $L$, all metric multiplication
operators are bounded and boundedly invertible.
Conjugating the actual weighted operator by
the identity isomorphism and using the displayed
trace-class factor gives
\[
 \begin{split}
 \operatorname{Tr}_{L^2(R_Ldr)}(M_\phi C_k^+Q_L^w)
 &=\operatorname{Tr}_{\mathcal K_+}
     \bigl(\mathfrak G_L^{-1}V^*M_{R_L}M_\phi VC_k^+\bigr).
 \end{split}
 \tag{MA-SB.18}
\]
Here $C_k^+V=VC_k^+$ because their Hardy
representatives are analytic multipliers.
No Fourier cutoff was moved past convolution.

Set $F_\phi=V^*M_\phi V$.
Since $\phi=1$ on the support of $R_L-1$,
\[
         R_L\phi=\phi+(R_L-1),\qquad
         V^*M_{R_L}M_\phi V
                           =F_\phi+\mathfrak G_L-1 .
 \tag{MA-SB.19}
\]
The operators $F_\phi C_k^+$ and
$(\mathfrak G_L-1)C_k^+$ are trace class by
the preceding compact-position argument.
Substitution in (MA-SB.18), followed only by
algebra of these trace-class products, gives
the exact trace correction
\[
 \begin{split}
 &\operatorname{Tr}_{L^2(R_Ldr)}(M_\phi C_k^+Q_L^w)
       -\operatorname{Tr}_{\mathcal K_+}
                                (M_\phi C_k^+VV^*)\\
 &\qquad=\operatorname{Tr}
             \bigl(D_L(F_\phi-1)C_k^+\bigr)
       =-\operatorname{Tr}
             \bigl(D_LV^*M_{1-\phi}VC_k^+\bigr).
 \end{split}
 \tag{MA-SB.20}
\]
In particular the right metric factor has not
been discarded from an uncompressed projection.
It was accounted for by (MA-SB.19) inside a
proved trace-class cyclic product.

\emph{Uniform control of the actual correction.}
The operator in the last trace of (MA-SB.20)
is trace class, as the difference of the
two preceding trace-class products.
Its range is in $P_L\mathcal K_+$.
Trace cyclicity with the bounded projection
$P_L$ therefore permits insertion of $P_L$
on its input:
\[
 \begin{split}
 \operatorname{Tr}
    (D_LV^*M_{1-\phi}VC_k^+)
 &=\operatorname{Tr}
    (D_LV^*M_{1-\phi}VC_k^+P_L),\\
 \left|\operatorname{Tr}
    (D_LV^*M_{1-\phi}VC_k^+)\right|
 &\le\|M_{1-\phi}VC_k^+P_L\|_1
 \le\nu_k(R-L).
 \end{split}
 \tag{MA-SB.21}
\]
The last inequality is exactly the translated
Hankel rectangle of MA-SB-001.
This estimates a scalar trace through an actual
trace-class factorization.  It does not assert
that an uncompressed weighted tail operator
has the same norm bound.

For the ordinary model term, $VV^*=1-P_K$,
and its identity contribution has trace zero.
The actual fixed lift $C_k^+P_K$ is trace class.
Hence
\[
 \begin{split}
 \operatorname{Tr}(M_\phi C_k^+VV^*)
 &=-\operatorname{Tr}(M_\phi C_k^+P_K),\\
 \left|\operatorname{Tr}(M_\phi C_k^+VV^*)+\tau_+(g)\right|
 &\le\omega_k(R).
 \end{split}
 \tag{MA-SB.22}
\]
Combining (MA-SB.20)--(MA-SB.22) proves
(MA-SB.13), with fixed tails independent of $L$.

\emph{Returning to the actual symmetric cutoff.}
The actual global PW range is supported on
$x\ge-\log Y$.  The convolution by $g$, supported
in $(0,\infty)$, preserves this half-line.
The global operator is zero on its orthogonal
complement and has range in it, so its trace
equals the trace of its restriction.
The actual unitary (MA-GP.84) converts that
restriction to (MA-SB.18), with
$\phi(r)=\rho(e^{r-L})$.

Under (MA-SB.14), every $r\ge0$ has
$e^{r-L}\ge1/Y\ge2/X>1/X+\delta$.
Thus the entire lower transition is absent
on this half-line.  Moreover $\phi(r)=1$
for
$0\le r\le L+\log(X-\delta)$.
Use this value of $R$ in (MA-SB.13).
It exceeds $L$, and $X-\delta\ge X/2$.
Both trace-norm tails are nonincreasing under
further left multiplication by tail projections,
so the second bound in (MA-SB.15) follows.
The eventual hypotheses in (MA-SB.14) hold
along every stated slow-band path.
Finally MA-GP-010 identifies the actual
finite-$S$ invariant global projector with
this same $Q_Y^{\rm glob}$, in the same norm.
\end{proof}

\subsection{The remaining correction is one-sided}

\paragraph{MA-SB-003 (actual slow transport and its precise geometric correction).}
The preceding theorem also applies to the common annular
profiles of MA-LA-014:
\[
       a_X(t)=A(t/X)-A(Xt),\qquad
       h(t)=g(\log t),
 \tag{MA-SB.23}
\]
where $A=1$ on $[0,1]$, $A=0$ on $[2,\infty)$,
and $0\le A\le1$ is smooth.
Indeed, for $X\ge2Y$ the lower term vanishes
on the actual global support half-line, and
the remaining profile is one through
$r=L+\log X$.
The same proof gives
\[
 \operatorname{Tr}_H(M_{a_X}C_gQ_Y^{\rm glob})
                         \longrightarrow-\tau_+(g)
       \qquad(X,Y\to\infty,\ Y/X\to0).
 \tag{MA-SB.24}
\]
All position multipliers in this display act through
$t=e^x$.

Choose $S$ to include every prime power norm in
the compact support of $h$; equivalently choose
it large enough for the support condition on
Masullo p.~520.  For the \emph{same} annular
position profile, the actual local quotient
operator of MA-LA-014 is
\[
 T^S_{X,Y,X^{-1}}(h)
       =M_{a_X}A_hD_SQ_Y^{\rm sm}D_S^*,
 \qquad
 Q_Y^{\rm sm}=\mathcal F_e^{-1}
                   M_{B_0(\,\cdot/Y)}\mathcal F_e ,
 \tag{MA-SB.25}
\]
with the fixed smooth even frequency profile
$B_0$ used there.
MA-LA-014 proves its actual trace limit
$A_0(g)+B_S(g)$ in this regime.
Since $B_S(g)=B(g)$ for the chosen support,
Meyer's actual common-space trace identity
MA-MEY identifies its polar-corrected limit
as $-\tau_{\rm all}(g)$, where
$\tau_{\rm all}$ is the trace of the full
smooth zero quotient.
Thus the limit of the difference of actual
traces is
\[
 \begin{split}
 &\lim_{\substack{X,Y\to\infty\\Y/X\to0}}
 \left\{\operatorname{Tr}_H(M_{a_X}C_gQ_Y^{\rm glob})
      -\left(\operatorname{Tr}T^S_{X,Y,X^{-1}}(h)-P(g)\right)\right\}\\
 &\hspace{35mm}=\tau_{\rm all}(g)-\tau_+(g).
 \end{split}
 \tag{MA-SB.26}
\]
The actual operators and their limits have been
constructed before this trace difference is
evaluated.  By the spectral identifications in
MA-GP-009 and Meyer, its value is
\[
                \sum_{\operatorname{Re}\rho\ge1/2+\alpha}
                        m_\rho\,\mathcal Mg(\rho-1/2).
 \tag{MA-SB.27}
\]
The upper-edge zeros are included, since the
half-plane Hardy quotient retains only the
strict inequality.
No such exterior or boundary zero is asserted
to exist.

For comparison, the actual completed full-strip
cutoff trace of MA-WP has retained zeros in
$1/2-\alpha<\operatorname{Re}\rho<1/2+\alpha$.
The slow-band limit minus that actual full-strip
trace is
\[
             -\sum_{\operatorname{Re}\rho\le1/2-\alpha}
                       m_\rho\,\mathcal Mg(\rho-1/2).
 \tag{MA-SB.28}
\]
This includes the lower edge and has the opposite
sign to the corresponding left-subset trace.
Both formulas concern an exact remaining
correction; neither establishes its nonvanishing.

The literal Gram-normalized finite cyclic operator
in MA-GP-008 supplies the same limit (MA-SB.24):
at each fixed $X,Y$ first use its proved
$\eta\downarrow0$ trace-norm limit, then the
uniform estimate (MA-SB.15).
Equivalently, along a specified slow-band sequence
choose $\eta_{X,Y}>0$ so that this fixed-data
Gram regularization changes the trace by at
most a prescribed error tending to zero.
This uses its two actual resolvents.
It asserts neither a uniform regularizer rate
nor that either Gram resolvent may be omitted.

\paragraph{Exact scope and trace-framework placement.}
The global conclusion (MA-SB.16) includes the source's
smooth symmetric thin-boundary cutoffs.
The paired local comparison (MA-SB.26) uses the
specified common fixed annular profiles for which
MA-LA-014 proves its local operator limit.
No unproved uniform estimate for an arbitrary
varying finite-place local smoothing is imported
by identifying those two profile conventions.

The polynomial weight is $s=0$ throughout the
application.  The abstract lemma MA-SB-001 applies
to any inner multiplier with the actual ambient
lift (MA-SB.2), but it does not alone identify a
positive-polynomial-weight Gram metric.
There is no generic-edge restriction, no
uniform-interior premise, and no use of
$\mathsf A$ or $M$.
The hypotheses of the actual half-plane range
theorem, including $0<\alpha<1/2$, remain explicit.

The finite-length zero modules determine the
quotient action, while the fixed trace-class
Hankel and model-space operators determine the
uniform tail control before scalar traces are
taken.  This is an actual trace-ideal comparison
in the inherited Grothendieck framework
(cote~1 scans~079--080 and~111--112,
BR-TR-01/02).  The new factorization is proved
above and uses no new archive citation.
It constructs no arithmetic positivity form,
Tate twist, or new zero-location theorem.
The source PW definitions and coherent global
Poisson normalization are precisely those
checked in MA-GP-003/010.
